\SourceLeaf{58}
\noindent Auch hier ist der Weg zum allgemeinen Fall schon gebahnt. Setzt man nämlich, wie es zulässig ist, die Ausdrücke gleich, die auf beiden Seiten mit $\sin(h+1)\varphi$ multipliziert sind, so ergibt sich

\[
\frac{\nu nA^{(h)}}{2}-\frac{\nu nA^{(h+2)}}{2}=(h+1)A^{(h+1)}+\frac{nhA^{(h)}}{2}+\frac{n(h+2)A^{(h+2)}}{2}.
\]

\noindent Löst man die Gleichung auf, so erhält man

\[
A^{(h+2)}=\frac{n(\nu-h)A^{(h)}-2(h+1)A^{(h+1)}}{n(\nu+h+2)},
\]

\noindent Diese allgemeine Formel enthält das Gesetz für jeden Koeffizienten der angesetzten Reihe, abgesehen von dem für $h=0$ erhaltenen $A^{(2)}$, der bereits durch eine gesonderte Untersuchung gefunden wurde. Sie liefert nämlich den jeweiligen Koeffizienten durch eine einfache Rechnung aus den beiden vorhergehenden. Zum Beispiel ergibt sich

\[
A^{(2)}=\frac{2(\nu nA-A^{(1)})}{n(\nu+2)};\qquad A^{(3)}=\frac{n(\nu-1)A^{(1)}-4A^{(2)}}{n(\nu+3)};
\]

\[
A^{(4)}=\frac{n(\nu-2)A^{(2)}-6A^{(3)}}{n(\nu+4)};\quad\mbox{etc.}
\]

$\delta$) Wenn aber $\nu$ eine negative Zahl ist, kann ein Nenner verschwinden und ein Koeffizient dadurch unbestimmt bleiben. Dann muss man auf die zuvor angegebenen unendlichen Reihen zurückgreifen. Dasselbe gilt für die beiden ersten Koeffizienten $A$ und $A^{(1)}$, die das Rekursionsschema bereits als bekannt voraussetzt. Mit einem besonderen Verfahren kann man jedoch $A^{(1)}$ durch $A$ ausdrücken. Es war nämlich

\[
A=1+.\,.\,.\,.+\frac{1}{2^{2h}}{}^{\nu}\mathfrak{B}^{(2h)}.{}^{2h}\mathfrak{B}^{(h)}n^{2h}+.\,.\,.,
\]

\[
A^{(1)}={}^{\nu}\mathfrak{B}^{(1)}n+.\,.\,.+\frac{1}{2^{2h}}{}^{\nu}\mathfrak{B}^{(2h+1)}.{}^{2h+1}\mathfrak{B}^{(h+1)}n^{2h+1}+.\,.\,..
\]

\noindent Nehmen wir $2nA+A^{(1)}$, so ist das Anfangsglied dieser Reihe

\[
=({}^{\nu}\mathfrak{B}^{(1)}+2)n=(\nu+2)n,
\]

\noindent und das allgemeine Glied $h^{\text{tus}}$

\[
2n . \frac{1}{2^{2h}}{}^{\nu}\mathfrak{B}^{(2h)}.{}^{2h}\mathfrak{B}^{(h)}n^{2h}+\frac{1}{2^{2h}}{}^{\nu}\mathfrak{B}^{(2h+1)}.{}^{2h+1}\mathfrak{B}^{(h+1)}n^{2h+1},
\]

\noindent Da

\[
{}^{\nu}\mathfrak{B}^{(2h+1)}={}^{\nu}\mathfrak{B}^{(2h)} . \frac{\nu-2h}{2h+1}\qquad\mbox{und}\qquad {}^{2h+1}\mathfrak{B}^{(h+1)}={}^{2h}\mathfrak{B}^{(h)} . \frac{2h+1}{h+1},
\]

\noindent gilt, vereinfacht sich dieses Glied zu

\[
\frac{1}{2^{2h}}{}^{\nu}\mathfrak{B}^{(2h)}.{}^{2h}\mathfrak{B}^{(h)}\left(2+\frac{\nu-2h}{2h+1} . \frac{2h+1}{h+1}\right)n^{2h+1}
\]

